\documentclass[11pt,a4paper]{article}
\usepackage[margin=27mm]{geometry}
\usepackage{amsmath,amssymb,amsthm,mathtools}
\usepackage{graphicx}
\usepackage{xcolor}
\usepackage{microtype}
\usepackage{enumitem}
\usepackage[colorlinks=true,linkcolor=blue!55!black,urlcolor=blue!55!black,citecolor=blue!55!black]{hyperref}

\newtheorem{theorem}{Theorem}
\newtheorem{lemma}{Lemma}
\newcommand{\R}{\mathbb R}
\newcommand{\e}{\mathrm e}
\newcommand{\thfour}{\theta_4}

\title{The scaled logarithmic derivative of $\theta_4$\\is decreasing and convex}
\author{Candidate full solution to an assertion in arXiv:1601.02972}
\date{11 August 2026}

\begin{document}
\maketitle

\begin{center}
\fcolorbox{blue!55!black}{blue!5}{\parbox{0.91\textwidth}{
\textbf{Status: candidate full result; likely valid; human review needed.}
For every $s>0$, the function
\[
 H(s)=s^2\frac{\thfour'(s)}{\thfour(s)}
\]
is strictly decreasing and strictly convex.}}
\end{center}

\section*{Source assertion}
Faulhuber and Steinerberger define
\[
 \thfour(s)=\sum_{k\in\mathbb Z}(-1)^k\e^{-\pi k^2s},\qquad s>0.
\]
After their refined logarithmic concavity theorem, they state that the stronger
claim that $s^2\thfour'(s)/\thfour(s)$ is decreasing and convex appears true,
but that their methods do not show it \cite[p.~3]{FS}.  The exact source text
is reproduced below.

\begin{center}
\includegraphics[width=0.96\textwidth]{figures/open_problem_crop.png}
\end{center}

\section*{Theorem and proof intuition}
\begin{theorem}\label{thm:main}
The function $H(s)=s^2\thfour'(s)/\thfour(s)$ satisfies
\[
 H'(s)<0\qquad\text{and}\qquad H''(s)>0
 \quad (s>0).
\]
\end{theorem}

\paragraph{Proof intuition.}
The two difficult ends of the parameter range want different expansions.
For $s\geq1$, Jacobi's product expresses $H$ as a positive sum of copies of
$x^2/(\e^x-1)$ evaluated at $x\geq\pi$; every copy is decreasing and convex.
For $s\leq1$, modular inversion exposes the leading behavior
$\pi/4-s/2$.  Its remaining terms can be paired as
$\e^{-C/s}-3\e^{-D/s}$ with $D>C$ and $D-C\geq2\pi$.
The exponential separation makes every pair increasing and convex.  Their
total slope is nevertheless less than $0.32$, so it cannot overcome the
leading slope $-1/2$; their positive curvatures establish convexity.

\section*{Two elementary estimates}
Set
\[
 \phi(x)=\frac{x^2}{\e^x-1}.
\]

\begin{lemma}\label{lem:phi}
$\phi'(x)<0$ and $\phi''(x)>0$ for every $x\geq\pi$.
\end{lemma}

\begin{proof}
Direct differentiation gives
\[
 \phi'(x)=\frac{x\bigl((2-x)\e^x-2\bigr)}{(\e^x-1)^2}<0
 \qquad(x\geq\pi>2).
\]
The numerator of $(\e^x-1)^3\phi''(x)$ is
\[
 N(x)=(x^2-4x+2)\e^{2x}+(x^2+4x-4)\e^x+2.
\]
Moreover
\[
 N'(x)=x\e^x\bigl((2x-6)\e^x+x+6\bigr)>0\qquad(x\geq\pi>3).
\]
It remains only to check the left endpoint without a decimal computation.
The elementary bounds $157/50<\pi<22/7$ and $\e^\pi<24$ imply, since
$\pi^2-4\pi+2<0$,
\begin{align*}
 \e^{-\pi}N(\pi)
 &>(24)(\pi^2-4\pi+2)+(\pi^2+4\pi-4)\\
 &=25\pi^2-92\pi+44
  >25(157/50)^2-92(157/50)+44>0.
\end{align*}
(For completeness, $\e^\pi<24$ follows from $\pi<3+1/7$,
$\e^3<20.1$, and $\e^{1/7}<7/6$.)  Thus $N(x)>0$ on the required range.
\end{proof}

\section*{The large-parameter half}
Jacobi's product is
\[
 \thfour(s)=\prod_{n\geq1}(1-\e^{-2\pi ns})
              (1-\e^{-(2n-1)\pi s})^2.
\]
Logarithmic differentiation yields
\begin{equation}\label{eq:large}
 H(s)=s^2\left[
 2\pi\sum_{n\geq1}\frac{n}{\e^{2\pi ns}-1}
 +2\pi\sum_{n\geq1}\frac{2n-1}{\e^{(2n-1)\pi s}-1}
 \right].
\end{equation}
Every summand in \eqref{eq:large} is a positive constant times
$\phi(c s)$, with $c\geq\pi$.  Hence Lemma~\ref{lem:phi} permits termwise
differentiation and gives
\begin{equation}\label{eq:large-signs}
 H'(s)<0,\qquad H''(s)>0\qquad(s\geq1).
\end{equation}

\section*{The small-parameter half}
The imaginary transformation and the product for $\theta_2$ give, for
$t>0$,
\[
 \thfour(s)=s^{-1/2}\theta_2(1/s),\qquad
 \theta_2(t)=2\e^{-\pi t/4}\prod_{n\geq1}
 (1-\e^{-2\pi nt})(1+\e^{-2\pi nt})^2.
\]
After logarithmic differentiation and expansion as geometric series,
\begin{equation}\label{eq:small}
 H(s)=\frac{\pi}{4}-\frac{s}{2}+R(s),
 \quad
 R(s)=2\pi\sum_{n\geq1}n\sum_{j\geq0}
 \left(\e^{-C_{n,j}/s}-3\e^{-D_{n,j}/s}\right),
\end{equation}
where
\[
 C_{n,j}=2\pi n(2j+1),\qquad D_{n,j}=2\pi n(2j+2).
\]
Indeed, if $a=2\pi n/s$, then the relevant bracket is
\[
 \frac{2}{\e^a+1}-\frac{1}{\e^a-1}
 =\frac{\e^a-3}{\e^{2a}-1}
 =\sum_{j\geq0}\left(\e^{-(2j+1)a}-3\e^{-(2j+2)a}\right).
\]

Fix $0<s\leq1$ and abbreviate $C=C_{n,j}$, $D=D_{n,j}$.
For $P(s)=\e^{-C/s}-3\e^{-D/s}$,
\begin{align}
 P'(s)&=\frac{C\e^{-C/s}}{s^2}
 \left[1-3\frac DC\e^{-(D-C)/s}\right],\label{eq:pprime}\\
 P''(s)&=\frac{C(C-2s)\e^{-C/s}}{s^4}
 \left[1-3\frac{D(D-2s)}{C(C-2s)}
 \e^{-(D-C)/s}\right].\label{eq:psecond}
\end{align}
Let $q=\e^{-2\pi}<1/500$.  Here $D/C\leq2$ and $D-C=2\pi n$,
so the bracket in \eqref{eq:pprime} is at least $1-6q>0$.
Also
\[
 \frac{D(D-2s)}{C(C-2s)}
 \leq2\left(2+\frac{2}{2\pi-2}\right)<5,
\]
so the bracket in \eqref{eq:psecond} is at least $1-15q>0$.
Consequently every pair is strictly increasing and strictly convex, and
therefore
\begin{equation}\label{eq:rconvex}
 R'(s)>0,\qquad R''(s)>0.
\end{equation}

It remains to keep the positive slope of $R$ below $1/2$.  Dropping the
negative member of every differentiated pair and observing that
$s^{-2}\e^{-C/s}$ increases on $(0,1]$, we obtain
\begin{align*}
 R'(s)
 &<4\pi^2\sum_{n\geq1}n^2\sum_{j\geq0}(2j+1)
 q^{n(2j+1)}\\
 &\leq4\pi^2\sum_{n\geq1}\frac{n^2q^n}{(1-q^n)^2}
 \leq4\pi^2\frac{q(1+q)}{(1-q)^5}<0.32<\frac12.
\end{align*}
In the last step we used $\pi^2<10$, $q<1/500$, $1+q<2$, and
$(1-q)^{-5}<2$.  Equations \eqref{eq:small} and \eqref{eq:rconvex} now give
\[
 H'(s)=-\frac12+R'(s)<0,
 \qquad H''(s)=R''(s)>0
 \qquad(0<s\leq1).
\]
Together with \eqref{eq:large-signs}, this proves
Theorem~\ref{thm:main}. \qed

\section*{Scope and novelty note}
This theorem proves only the precise theta-function assertion on source PDF
page 3.  It does not resolve the paper's much broader Strohmer--Beaver
hexagonal-lattice conjecture.  The run's four cheap indexes and targeted arXiv
searches for the exact scaled logarithmic derivative statement found the
source paper but no later primary-source resolution.  Novelty therefore
remains provisional pending expert literature review.

\begin{thebibliography}{9}
\bibitem{FS}
M.~Faulhuber and S.~Steinerberger,
\emph{Optimal Gabor frame bounds for separable lattices and estimates for
Jacobi theta functions}, arXiv:1601.02972v3 (2016).
\end{thebibliography}

\end{document}

